% ===== Préambule commun aux documents Séries Chronologiques =====
\usepackage[utf8]{inputenc}
\usepackage[T1]{fontenc}
\usepackage[french]{babel}
\usepackage{lmodern}
\usepackage{amsmath,amssymb,amsthm,mathtools}
\usepackage{geometry}
\geometry{a4paper,margin=2.3cm}
\usepackage{enumitem}
\usepackage{array,booktabs,longtable}
\usepackage{xcolor}
\usepackage[most]{tcolorbox}
\usepackage{etoolbox}
\usepackage{graphicx}
\graphicspath{{figs/}}
\usepackage{caption}
\captionsetup{font=small,labelfont={bf,color=insea}}
\usepackage{fancyhdr}
\usepackage{titlesec}
\usepackage{hyperref}
\hypersetup{colorlinks=true,linkcolor=blue!55!black,urlcolor=blue!55!black,
  citecolor=blue!55!black,bookmarksopen=true}
\usepackage{bookmark}

% --- Couleurs ---
\definecolor{insea}{RGB}{0,76,128}
\definecolor{insealight}{RGB}{232,240,247}
\definecolor{vert}{RGB}{0,120,70}
\definecolor{vertlight}{RGB}{232,245,238}
\definecolor{orange}{RGB}{190,90,0}
\definecolor{orangelight}{RGB}{252,242,228}
\definecolor{gris}{RGB}{90,90,90}

% --- En-têtes / pieds ---
\pagestyle{fancy}
\fancyhf{}
\fancyhead[L]{\small\color{gris}\leftmark}
\fancyhead[R]{\small\color{gris}Séries Chronologiques}
\fancyfoot[C]{\thepage}
\renewcommand{\headrulewidth}{0.4pt}
\renewcommand{\footrulewidth}{0pt}

% --- Titres de section colorés ---
\titleformat{\section}{\Large\bfseries\color{insea}}{\thesection.}{0.6em}{}
\titleformat{\subsection}{\large\bfseries\color{insea!85!black}}{\thesubsection.}{0.6em}{}
\titleformat{\subsubsection}{\bfseries\color{gris}}{\thesubsubsection.}{0.6em}{}

% --- Théorèmes (pour le cours) ---
\theoremstyle{definition}
\newtheorem{definition}{Définition}[section]
\newtheorem{exemple}{Exemple}[section]
\newtheorem{remarque}{Remarque}[section]
\newtheorem{methode}{Méthode}[section]
\theoremstyle{plain}
\newtheorem{theoreme}{Théorème}[section]
\newtheorem{proposition}{Proposition}[section]
\newtheorem{corollaire}{Corollaire}[section]
\newtheorem{propriete}{Propriété}[section]

% --- Boîtes ---
\newtcolorbox{recette}[1][]{colback=vertlight,colframe=vert,
  fonttitle=\bfseries,title={Recette d'examen\ifblank{#1}{}{ --- #1}}}
\newtcolorbox{attention}[1][]{colback=orangelight,colframe=orange,
  fonttitle=\bfseries,title={Attention\ifblank{#1}{}{ --- #1}}}
\newtcolorbox{rappel}[1][]{colback=insealight,colframe=insea,
  fonttitle=\bfseries,title={Rappel de cours\ifblank{#1}{}{ --- #1}}}
\newtcolorbox{reponse}[1][]{colback=vertlight,colframe=vert!70!black,
  left=2mm,right=2mm,top=1mm,bottom=1mm,boxrule=0.8pt,#1}

% --- Notations Séries Chronologiques ---
\newcommand{\E}{\mathbb{E}}
\newcommand{\Var}{\operatorname{Var}}
\newcommand{\Cov}{\operatorname{Cov}}
\newcommand{\Corr}{\operatorname{Corr}}
\newcommand{\cor}{\operatorname{cor}}
\newcommand{\BB}{\operatorname{BB}}
\newcommand{\R}{\mathbb{R}}
\newcommand{\Z}{\mathbb{Z}}
\newcommand{\N}{\mathbb{N}}
\newcommand{\Xhat}{\widehat{X}}
\newcommand{\phihat}{\widehat{\phi}}
\newcommand{\rhohat}{\widehat{\rho}}
\newcommand{\gammahat}{\widehat{\gamma}}
\newcommand{\eps}{\varepsilon}
\DeclareMathOperator{\sgn}{sgn}

% Étiquette « bonne réponse » pour les QCM (utilisée dans le corrigé)
\newcommand{\bonne}[1]{\textbf{\color{vert}#1}}
\newcommand{\horsprog}{\textsuperscript{\color{orange}\,$\star$}}
